% from main.tex
%%%%%%%%%%  scan idx 611  |  folio 598  %%%%%%%%%%
Donc \(P\to P/R\), qui fait de \(P\) un \(R\)-torseur sur \(P/R\), est affine. D'ailleurs on sait (SGA 3 1.4.2) que \(P/R\to Q\) est une immersion fermée, donc affine, donc il en est de même du composé \(P\to P/R\to Q\).

\textbf{Corollaire 1.5.} Soit \(k\) un corps.
\begin{enumerate}[label=\alph*)]
\item Soit \(X\) un \(k\)-schéma propre. Alors \(\underline{\Pic}_{X/k}\) est représentable (J. P. Murre).
\item Soit \(f:X\to Y\) un morphisme surjectif entre deux \(k\)-schémas propres. Alors le morphisme induit \(f^*:\underline{\Pic}_{Y/k}\longrightarrow \underline{\Pic}_{X/k}\) est affine.
\end{enumerate}

\textbf{Corollaire 1.6.} Soient \(S\) un schéma noethérien intègre, \(f:X\to Y\) un \(S\)-morphisme surjectif entre deux \(S\)-schémas propres.
\begin{enumerate}[label=\alph*)]
\item Si \(Y\) est normal sur \(S\) (c'est-à-dire plat à fibres géométriquement normales), alors, quitte à restreindre \(S\), \(f^*:\underline{\Pic}_{Y/S}\to\underline{\Pic}_{X/S}\) est quasi-fini.
\item Si \(\underline{\Pic}_{Y/S}\) est essentiellement propre sur \(S\) (c'est-à-dire, s'il satisfait au critère voluatif de propreté, par exemple, si \(Y\) est lisse sur \(S\)) alors quitte à restreindre \(S\), \(f^*:\underline{\Pic}_{Y/S}\to\underline{\Pic}_{X/S}\) est fini (et en particulier, affine).
\end{enumerate}

En effet dans a), pour tout \(s\in S\), \(R_s=\Ker(f_s^*)\) est un groupe algébrique affine~; sa composante neutre est un groupe algébrique propre sur \(k\) (FGA 236, Théorème 2.1 (iii))~; \(R_s\) est donc fini. Dans b) \(f^*\) est à la fois propre et quasi-affine, donc fini.

\textbf{Remarque 1.7.} On ignore si 1.1 reste valable quand on y remplace « quasi-affine » par affine, comme pourraient le suggérer 1.5 et 1.6.

